\documentclass[12pt,twoside,openany]{memoir}
%\usepackage{mlmodern}
%\usepackage{tgpagella} % text only
%\usepackage{mathpazo}  % math & text
\usepackage[T1]{fontenc}
\usepackage{amsmath}
\usepackage[fixamsmath]{mathtools}  % Extension to amsmath
\usepackage{amsthm}
\usepackage{amssymb}
\renewcommand*{\mathbf}[1]{\varmathbb{#1}}
\usepackage{newpxtext}
\usepackage{eulerpx}
\usepackage{mathrsfs}

\usepackage{eucal}
\usepackage{datetime}
    \newdateformat{specialdate}{\THEYEAR\ \monthname\ \THEDAY}
\usepackage[margin=1in]{geometry}
\usepackage{fancyhdr}
    \pagestyle{fancy}
    \renewcommand{\headrulewidth}{0pt}
    \cfoot{\scriptsize \thepage}
\usepackage{keytheorems}
    \newkeytheorem{theorem}[
    name = Theorem,
    numberwithin = section,
    ]
    \newkeytheorem{theorem*}[
    name=Theorem,
    numbered=false
    ]
    \newkeytheorem{proposition}[
    name = Proposition,
    numberwithin = section,
    sibling=theorem
    ]
    \newkeytheorem{lemma}[
    name = Lemma,
    numberwithin = section,
    sibling=theorem
    ]
    \newkeytheorem{corollary}[
    name = Corollary,
    numberwithin = section,
    sibling=theorem
    ]
    \newkeytheorem{definition}[
    name = Definition,
    numberwithin = section,
    style = definition
    ]
    \newkeytheorem{example}[
    name = Example,
    numberwithin = section,
    style = definition
    ]
    \newkeytheorem{remark}[
    name=Remark,
    style=remark, 
    numbered=false  
    ]
    \newkeytheorem{answer}[
    name=Answer,
    style=remark, 
    numbered=false  
    ]
    \newkeytheorem{solution}[
    name=Solution,
    style=remark, 
    numbered=false  
    ]
    \newkeytheorem{exercise}[
    name = Exercise
    ]
    \newkeytheorem{axiom}[name = Axiom]
\usepackage{enumitem}
\usepackage{titlesec}
    \titleformat{\chapter}[display]
    {\bfseries\LARGE\raggedright}
    {Chapter {\thechapter}}
    {.1ex minus .1ex}
    {\Large}
    \titlespacing{\chapter}
    {3pc}{*3}{40pt}[3pc]

    \titleformat{\section}[block]
    {\normalfont\bfseries\large}
    {\S\ \thesection.}{.5em}{}[]
    \titlespacing{\section}
    {0pt}{3ex plus .1ex minus .2ex}{3ex plus .1ex minus .2ex}
\usepackage[utf8x]{inputenc}
\usepackage{tikz}
\usepackage{tikz-cd}
\usepackage{float}
\usetikzlibrary{patterns}
\usetikzlibrary{positioning,arrows.meta}
\usetikzlibrary{decorations.pathreplacing}
\usepackage{wasysym}
\usepackage{hyperref}
\usepackage{ulem}
\usepackage{pgfplots}
\usepackage{xcolor}
\usepackage{matlab-prettifier}
\lstdefinestyle{Matlab-console}{
  basicstyle=\ttfamily\small,
  numbers=none,
  frame=single,
  breaklines=true,
  showstringspaces=false 
}
\definecolor{matlabbg}{HTML}{F2F2F2}
\definecolor{matlabln}{HTML}{6E6E6E}

\lstdefinestyle{MatlabEditorGrey}{
  style=Matlab-editor,
  backgroundcolor=\color{matlabbg},
  numbers=left,
  numberstyle=\color{matlabln}\ttfamily\scriptsize,
  stepnumber=1,
  numbersep=10pt,
  xleftmargin=2.2em,     % space for line numbers
  frame=single,
  framerule=0pt,         % no border line (keep the box spacing)
  framesep=6pt,
  breaklines=true
}
\usetikzlibrary{patterns}
\usepgfplotslibrary{fillbetween}
\pgfplotsset{compat=1.18}
\linespread{1.1}
%%%%%%%%%%%%%%%%%%%%%%%%%%%%%%%%%%%%%%%%%%%%%%%%%%%%%%%%%%%%%
%%%%%%%%%%%%%%%%%%%%%%%%%%%%%%%%%%%%%%%%%%%%%%%%%%%%%%%%%%%%%
% ============================================================
% makros.tex — reorganized, full restored version
% ============================================================

% ============================================================
% Sha / Cyrillic support notes
% ============================================================

%to make the correct symbol for Sha
%\newcommand\cyr{%
%\renewcommand\rmdefault{wncyr}%
%\renewcommand\sfdefault{wncyss}%
%\renewcommand\encodingdefault{OT2}%
%\normalfont \selectfont} \DeclareTextFontCommand{\textcyr}{\cyr}

% ============================================================
% Math operators
% ============================================================

\DeclareMathOperator{\ab}{ab}
\DeclareMathOperator{\ad}{ad}
\DeclareMathOperator{\adj}{adj}
\DeclareMathOperator{\alg}{alg}
\DeclareMathOperator{\Alt}{Alt}
\DeclareMathOperator{\Ann}{Ann}
\DeclareMathOperator{\arith}{arith}
\DeclareMathOperator{\Aut}{Aut}
\DeclareMathOperator{\Be}{B}
\DeclareMathOperator{\Bd}{Bd}
\DeclareMathOperator{\card}{card}
\DeclareMathOperator{\Char}{char}
\DeclareMathOperator{\csp}{csp}
\DeclareMathOperator{\codim}{codim}
\DeclareMathOperator{\coker}{coker}
\DeclareMathOperator{\coh}{H}
\DeclareMathOperator{\compl}{compl}
\DeclareMathOperator{\conj}{conj}
\DeclareMathOperator{\cont}{cont}
\DeclareMathOperator{\Cov}{Cov}
\DeclareMathOperator{\crys}{crys}
\DeclareMathOperator{\Crys}{Crys}
\DeclareMathOperator{\cusp}{cusp}
\DeclareMathOperator{\diag}{diag}
\DeclareMathOperator{\diam}{diam}
\DeclareMathOperator{\Dom}{Dom}
\DeclareMathOperator{\disc}{disc}
\DeclareMathOperator{\dist}{dist}
\DeclareMathOperator{\Div}{div}
\DeclareMathOperator{\dR}{dR}
\DeclareMathOperator{\Eis}{Eis}
\DeclareMathOperator{\End}{End}
\DeclareMathOperator{\ev}{ev}
\DeclareMathOperator{\eval}{eval}
\DeclareMathOperator{\Eq}{Eq}
\DeclareMathOperator{\Ext}{Ext}
\DeclareMathOperator{\Fil}{Fil}
\DeclareMathOperator{\Fitt}{Fitt}
\DeclareMathOperator{\Frob}{Frob}
\DeclareMathOperator{\G}{G}
\DeclareMathOperator{\Gal}{Gal}
\DeclareMathOperator{\GL}{GL}
\DeclareMathOperator{\Gr}{Gr}
\DeclareMathOperator{\Graph}{Graph}
\DeclareMathOperator{\GSp}{GSp}
\DeclareMathOperator{\GUn}{GU}
\DeclareMathOperator{\Hom}{Hom}
\DeclareMathOperator{\id}{id}
\DeclareMathOperator{\Id}{Id}
\DeclareMathOperator{\Ik}{Ik}
\DeclareMathOperator{\IM}{Im}
\DeclareMathOperator{\Image}{im}
\DeclareMathOperator{\Ind}{Ind}
\DeclareMathOperator{\Inf}{inf}
\DeclareMathOperator{\ins}{ins}
\DeclareMathOperator{\inv}{inv}
\DeclareMathOperator{\Isom}{Isom}
\DeclareMathOperator{\J}{J}
\DeclareMathOperator{\Jac}{Jac}
\DeclareMathOperator{\lcm}{lcm}
\DeclareMathOperator{\length}{length}
\DeclareMathOperator*{\limit}{limit}
\DeclareMathOperator{\Log}{Log}
\DeclareMathOperator{\M}{M}
\DeclareMathOperator{\Mat}{Mat}
\DeclareMathOperator{\N}{N}
\DeclareMathOperator{\Nm}{Nm}
\DeclareMathOperator{\NIk}{N-Ik}
\DeclareMathOperator{\NSK}{N-SK}
\DeclareMathOperator{\new}{new}
\DeclareMathOperator{\obj}{obj}
\DeclareMathOperator{\old}{old}
\DeclareMathOperator{\ord}{ord}
\DeclareMathOperator{\Or}{O}
\DeclareMathOperator{\op}{op}
\DeclareMathOperator{\out}{out}
\DeclareMathOperator{\PGL}{PGL}
\DeclareMathOperator{\PGSp}{PGSp}
\DeclareMathOperator{\rank}{rank}
\DeclareMathOperator{\Ran}{Ran}
\DeclareMathOperator{\rel}{Rel}
\DeclareMathOperator{\Real}{Re}
\DeclareMathOperator{\RES}{res}
\DeclareMathOperator{\Res}{Res}
%\DeclareMathOperator{\Sha}{\textcyr{Sh}}
\DeclareMathOperator{\Sel}{Sel}
\DeclareMathOperator{\semi}{ss}
\DeclareMathOperator{\sgn}{sign}
\DeclareMathOperator{\SK}{SK}
\DeclareMathOperator{\SL}{SL}
\DeclareMathOperator{\SO}{SO}
\DeclareMathOperator{\Sp}{Sp}
\DeclareMathOperator{\Span}{span}
\DeclareMathOperator{\Spec}{Spec}
\DeclareMathOperator{\spin}{spin}
\DeclareMathOperator{\st}{st}
\DeclareMathOperator{\St}{St}
\DeclareMathOperator{\SUn}{SU}
\DeclareMathOperator{\supp}{supp}
\DeclareMathOperator{\Sup}{sup}
\DeclareMathOperator{\Sym}{Sym}
\DeclareMathOperator{\Tam}{Tam}
\DeclareMathOperator{\tors}{tors}
\DeclareMathOperator{\tr}{tr}
\DeclareMathOperator{\Tr}{Tr}
\DeclareMathOperator{\un}{un}
\DeclareMathOperator{\Un}{U}
\DeclareMathOperator{\val}{val}
\DeclareMathOperator{\vol}{vol}

% ============================================================
% General aliases and short macros
% ============================================================

\newcommand{\absgal}{\G_{\bbQ}}
\newcommand{\Loc}{L_{\operatorname{loc}}}
\renewcommand{\k}{\kappa}
\newcommand{\Ff}{F_{f}}
%\newcommand{\ts}{\,^{t}\!}
\newcommand{\lat}{\mathscr{L}}
\newcommand{\ov}[1]{\overline{#1}}
\newcommand{\up}{\upsilon}
\newcommand{\e}{\epsilon}
\newcommand{\tac}{\textasteriskcentered}

%mahesh macros
\newcommand{\tm}{\textrm}

\newcommand{\bone}{\mathbf{1}}
\newcommand{\sign}{\mathrm{sign}}
\newcommand{\eps}{\varepsilon}
\newcommand{\textui}[1]{\underline{\textit{#1}}}

%\newcommand{\ov}{\overline}
%\newcommand{\un}{\underline}
\newcommand{\fin}{\mathrm{fin}}
\newcommand{\nl}{\newline \mbox{}}
\newcommand{\h}[1]{\hspace{#1pt}}
\newcommand{\mtext}[1]{\hspace{6pt}\text{#1}\hspace{6pt}}
\newcommand{\lnorm}{\left\lVert}
\newcommand{\rnorm}{\right\rVert}
\newcommand{\ds}{\displaystyle}
\newcommand{\ts}{\textstyle}
\newcommand{\sfrac}[2]{{}^{#1}\mskip -5mu/\mskip -3mu_{#2}}

% ============================================================
% Category-style operators and contoured variants
% ============================================================

\DeclareMathOperator{\Sets}{S \mkern1.04mu e \mkern1.04mu t \mkern1.04mu s}
    \newcommand{\cSets}{\scalebox{1.02}{\contour{black}{$\Sets$}}}
    
\DeclareMathOperator{\Groups}{G \mkern1.04mu r \mkern1.04mu o \mkern1.04mu u \mkern1.04mu p \mkern1.04mu s}
    \newcommand{\cGroups}{\scalebox{1.02}{\contour{black}{$\Groups$}}}

\DeclareMathOperator{\TTop}{T \mkern1.04mu o \mkern1.04mu p}
    \newcommand{\cTop}{\scalebox{1.02}{\contour{black}{$\TTop$}}}

\DeclareMathOperator{\Htp}{H \mkern1.04mu t \mkern1.04mu p}
    \newcommand{\cHtp}{\scalebox{1.02}{\contour{black}{$\Htp$}}}

\DeclareMathOperator{\Mod}{M \mkern1.04mu o \mkern1.04mu d}
    \newcommand{\cMod}{\scalebox{1.02}{\contour{black}{$\Mod$}}}

\DeclareMathOperator{\Ab}{A \mkern1.04mu b}
    \newcommand{\cAb}{\scalebox{1.02}{\contour{black}{$\Ab$}}}

\DeclareMathOperator{\Rings}{R \mkern1.04mu i \mkern1.04mu n \mkern1.04mu g \mkern1.04mu s}
    \newcommand{\cRings}{\scalebox{1.02}{\contour{black}{$\Rings$}}}

\DeclareMathOperator{\ComRings}{C \mkern1.04mu o \mkern1.04mu m \mkern1.04mu R \mkern1.04mu i \mkern1.04mu n \mkern1.04mu g \mkern1.04mu s}
    \newcommand{\cComRings}{\scalebox{1.05}{\contour{black}{$\ComRings$}}}

\DeclareMathOperator{\hHom}{H \mkern1.04mu o \mkern1.04mu m}
    \newcommand{\cHom}{\scalebox{1.02}{\contour{black}{$\hHom$}}}

% ============================================================
% Mathcal
% ============================================================

\newcommand{\cA}{\mathcal{A}}
\newcommand{\cB}{\mathcal{B}}
\newcommand{\cC}{\mathcal{C}}
\newcommand{\cD}{\mathcal{D}}
\newcommand{\cE}{\mathcal{E}}
\newcommand{\cF}{\mathcal{F}}
\newcommand{\cG}{\mathcal{G}}
\newcommand{\cH}{\mathcal{H}}
\newcommand{\cI}{\mathcal{I}}
\newcommand{\cJ}{\mathcal{J}}
\newcommand{\cK}{\mathcal{K}}
\newcommand{\cL}{\mathcal{L}}
\newcommand{\cM}{\mathcal{M}}
\newcommand{\cN}{\mathcal{N}}
\newcommand{\cO}{\mathcal{O}}
\newcommand{\cP}{\mathcal{P}}
\newcommand{\cQ}{\mathcal{Q}}
\newcommand{\cR}{\mathcal{R}}
\newcommand{\cS}{\mathcal{S}}
\newcommand{\cT}{\mathcal{T}}
\newcommand{\cU}{\mathcal{U}}
\newcommand{\cV}{\mathcal{V}}
\newcommand{\cW}{\mathcal{W}}
\newcommand{\cX}{\mathcal{X}}
\newcommand{\cY}{\mathcal{Y}}
\newcommand{\cZ}{\mathcal{Z}}

% ============================================================
% Mathscrl% 
%============================================================

\newcommand{\scra}{\mathscr{a}}
\newcommand{\scrA}{\mathscr{A}}
\newcommand{\scrb}{\mathscr{b}}
\newcommand{\scrB}{\mathscr{B}}
\newcommand{\scrc}{\mathscr{c}}
\newcommand{\scrC}{\mathscr{C}}
\newcommand{\scrd}{\mathscr{d}}
\newcommand{\scrD}{\mathscr{D}}
\newcommand{\scre}{\mathscr{e}}
\newcommand{\scrE}{\mathscr{E}}
\newcommand{\scrf}{\mathscr{f}}
\newcommand{\scrF}{\mathscr{F}}
\newcommand{\scrg}{\mathscr{g}}
\newcommand{\scrG}{\mathscr{G}}
\newcommand{\scrh}{\mathscr{h}}
\newcommand{\scrH}{\mathscr{H}}
\newcommand{\scri}{\mathscr{i}}
\newcommand{\scrI}{\mathscr{I}}
\newcommand{\scrj}{\mathscr{j}}
\newcommand{\scrJ}{\mathscr{J}}
\newcommand{\scrk}{\mathscr{k}}
\newcommand{\scrK}{\mathscr{K}}
\newcommand{\scrl}{\mathscr{l}}
\newcommand{\scrL}{\mathscr{L}}
\newcommand{\scrm}{\mathscr{m}}
\newcommand{\scrM}{\mathscr{M}}
\newcommand{\scrn}{\mathscr{n}}
\newcommand{\scrN}{\mathscr{N}}
\newcommand{\scro}{\mathscr{o}}
\newcommand{\scrO}{\mathscr{O}}
\newcommand{\scrp}{\mathscr{p}}
\newcommand{\scrP}{\mathscr{P}}
\newcommand{\scrq}{\mathscr{q}}
\newcommand{\scrQ}{\mathscr{Q}}
\newcommand{\scrr}{\mathscr{r}}
\newcommand{\scrR}{\mathscr{R}}
\newcommand{\scrs}{\mathscr{s}}
\newcommand{\scrS}{\mathscr{S}}
\newcommand{\scrt}{\mathscr{t}}
\newcommand{\scrT}{\mathscr{T}}
\newcommand{\scru}{\mathscr{u}}
\newcommand{\scrU}{\mathscr{U}}
\newcommand{\scrv}{\mathscr{v}}
\newcommand{\scrV}{\mathscr{V}}
\newcommand{\scrw}{\mathscr{w}}
\newcommand{\scrW}{\mathscr{W}}
\newcommand{\scrx}{\mathscr{x}}
\newcommand{\scrX}{\mathscr{X}}
\newcommand{\scry}{\mathscr{y}}
\newcommand{\scrY}{\mathscr{Y}}
\newcommand{\scrz}{\mathscr{z}}
\newcommand{\scrZ}{\mathscr{Z}}

% ============================================================
% Mathfrak (missing \fi)
% ============================================================

\newcommand{\fa}{\mathfrak{a}}
\newcommand{\fA}{\mathfrak{A}}
\newcommand{\fb}{\mathfrak{b}}
\newcommand{\fB}{\mathfrak{B}}
\newcommand{\fc}{\mathfrak{c}}
\newcommand{\fC}{\mathfrak{C}}
\newcommand{\fd}{\mathfrak{d}}
\newcommand{\fD}{\mathfrak{D}}
\newcommand{\fe}{\mathfrak{e}}
\newcommand{\fE}{\mathfrak{E}}
\newcommand{\ff}{\mathfrak{f}}
\newcommand{\fF}{\mathfrak{F}}
\newcommand{\fg}{\mathfrak{g}}
\newcommand{\fG}{\mathfrak{G}}
\newcommand{\fh}{\mathfrak{h}}
\newcommand{\fH}{\mathfrak{H}}
\newcommand{\fI}{\mathfrak{I}}
\newcommand{\fj}{\mathfrak{j}}
\newcommand{\fJ}{\mathfrak{J}}
\newcommand{\fk}{\mathfrak{k}}
\newcommand{\fK}{\mathfrak{K}}
\newcommand{\fl}{\mathfrak{l}}
\newcommand{\fL}{\mathfrak{L}}
\newcommand{\fm}{\mathfrak{m}}
\newcommand{\fM}{\mathfrak{M}}
\newcommand{\fn}{\mathfrak{n}}
\newcommand{\fN}{\mathfrak{N}}
\newcommand{\fo}{\mathfrak{o}}
\newcommand{\fO}{\mathfrak{O}}
\newcommand{\fp}{\mathfrak{p}}
\newcommand{\fP}{\mathfrak{P}}
\newcommand{\fq}{\mathfrak{q}}
\newcommand{\fQ}{\mathfrak{Q}}
\newcommand{\fr}{\mathfrak{r}}
\newcommand{\fR}{\mathfrak{R}}
\newcommand{\fs}{\mathfrak{s}}
\newcommand{\fS}{\mathfrak{S}}
\newcommand{\ft}{\mathfrak{t}}
\newcommand{\fT}{\mathfrak{T}}
\newcommand{\fu}{\mathfrak{u}}
\newcommand{\fU}{\mathfrak{U}}
\newcommand{\fv}{\mathfrak{v}}
\newcommand{\fV}{\mathfrak{V}}
\newcommand{\fw}{\mathfrak{w}}
\newcommand{\fW}{\mathfrak{W}}
\newcommand{\fx}{\mathfrak{x}}
\newcommand{\fX}{\mathfrak{X}}
\newcommand{\fy}{\mathfrak{y}}
\newcommand{\fY}{\mathfrak{Y}}
\newcommand{\fz}{\mathfrak{z}}
\newcommand{\fZ}{\mathfrak{Z}}

% ============================================================
% Mathbf
% ============================================================

\newcommand{\bfA}{\mathbf{A}}
\newcommand{\bfB}{\mathbf{B}}
\newcommand{\bfC}{\mathbf{C}}
\newcommand{\bfD}{\mathbf{D}}
\newcommand{\bfE}{\mathbf{E}}
\newcommand{\bfF}{\mathbf{F}}
\newcommand{\bfG}{\mathbf{G}}
\newcommand{\bfH}{\mathbf{H}}
\newcommand{\bfI}{\mathbf{I}}
\newcommand{\bfJ}{\mathbf{J}}
\newcommand{\bfK}{\mathbf{K}}
\newcommand{\bfL}{\mathbf{L}}
\newcommand{\bfM}{\mathbf{M}}
\newcommand{\bfN}{\mathbf{N}}
\newcommand{\bfO}{\mathbf{O}}
\newcommand{\bfP}{\mathbf{P}}
\newcommand{\bfQ}{\mathbf{Q}}
\newcommand{\bfR}{\mathbf{R}}
\newcommand{\bfS}{\mathbf{S}}
\newcommand{\bfT}{\mathbf{T}}
\newcommand{\bfU}{\mathbf{U}}
\newcommand{\bfV}{\mathbf{V}}
\newcommand{\bfW}{\mathbf{W}}
\newcommand{\bfX}{\mathbf{X}}
\newcommand{\bfY}{\mathbf{Y}}
\newcommand{\bfZ}{\mathbf{Z}}

\newcommand{\bfa}{\mathbf{a}}
\newcommand{\bfb}{\mathbf{b}}
\newcommand{\bfc}{\mathbf{c}}
\newcommand{\bfd}{\mathbf{d}}
\newcommand{\bfe}{\mathbf{e}}
\newcommand{\bff}{\mathbf{f}}
\newcommand{\bfg}{\mathbf{g}}
\newcommand{\bfh}{\mathbf{h}}
\newcommand{\bfi}{\mathbf{i}}
\newcommand{\bfj}{\mathbf{j}}
\newcommand{\bfk}{\mathbf{k}}
\newcommand{\bfl}{\mathbf{l}}
\newcommand{\bfm}{\mathbf{m}}
\newcommand{\bfn}{\mathbf{n}}
\newcommand{\bfo}{\mathbf{o}}
\newcommand{\bfp}{\mathbf{p}}
\newcommand{\bfq}{\mathbf{q}}
\newcommand{\bfr}{\mathbf{r}}
\newcommand{\bfs}{\mathbf{s}}
\newcommand{\bft}{\mathbf{t}}
\newcommand{\bfu}{\mathbf{u}}
\newcommand{\bfv}{\mathbf{v}}
\newcommand{\bfw}{\mathbf{w}}
\newcommand{\bfx}{\mathbf{x}}
\newcommand{\bfy}{\mathbf{y}}
\newcommand{\bfz}{\mathbf{z}}

% ============================================================
% Blackboard bold
% ============================================================

\newcommand{\bbA}{\mathbb{A}}
\newcommand{\bbB}{\mathbb{B}}
\newcommand{\bbC}{\mathbb{C}}
\newcommand{\bbD}{\mathbb{D}}
\newcommand{\bbE}{\mathbb{E}}
\newcommand{\bbF}{\mathbb{F}}
\newcommand{\bbG}{\mathbb{G}}
\newcommand{\bbH}{\mathbb{H}}
\newcommand{\bbI}{\mathbb{I}}
\newcommand{\bbJ}{\mathbb{J}}
\newcommand{\bbK}{\mathbb{K}}
\newcommand{\bbL}{\mathbb{L}}
\newcommand{\bbM}{\mathbb{M}}
\newcommand{\bbN}{\mathbb{N}}
\newcommand{\bbO}{\mathbb{O}}
\newcommand{\bbP}{\mathbb{P}}
\newcommand{\bbQ}{\mathbb{Q}}
\newcommand{\bbR}{\mathbb{R}}
\newcommand{\bbS}{\mathbb{S}}
\newcommand{\bbT}{\mathbb{T}}
\newcommand{\bbU}{\mathbb{U}}
\newcommand{\bbV}{\mathbb{V}}
\newcommand{\bbW}{\mathbb{W}}
\newcommand{\bbX}{\mathbb{X}}
\newcommand{\bbY}{\mathbb{Y}}
\newcommand{\bbZ}{\mathbb{Z}}

\newcommand{\Bmu}{\mbox{$\raisebox{-0.59ex}
  {$l$}\hspace{-0.18em}\mu\hspace{-0.88em}\raisebox{-0.98ex}{\scalebox{2}
  {$\color{white}.$}}\hspace{-0.416em}\raisebox{+0.88ex}
  {$\color{white}.$}\hspace{0.46em}$}{}}  %need graphicx and xcolor. this produces blackboard bold mu 

% ============================================================
% Greek-letter aliases
% ============================================================

\newcommand{\jota}{\jmath}

% ============================================================
% Matrices
% ============================================================

\newcommand{\bmat}{\left( \begin{matrix}}
\newcommand{\emat}{\end{matrix} \right)}

\newcommand{\bbmat}{\left[ \begin{matrix}}
\newcommand{\ebmat}{\end{matrix} \right]}

\newcommand{\pmat}{\left( \begin{smallmatrix}}
\newcommand{\epmat}{\end{smallmatrix} \right)}

\newcommand{\mat}[4]{\begin{pmatrix}{#1}&{#2}\\{#3}&{#4}\end{pmatrix}}

% ============================================================
% Residues and averaged integrals
% ============================================================

\newcommand{\res}[1]{\underset{#1}{\RES}\,}

\makeatletter
\newcommand*{\mint}[1]{%
  % #1: overlay symbol
  \mint@l{#1}{}%
}
\newcommand*{\mint@l}[2]{%
  % #1: overlay symbol
  % #2: limits
  \@ifnextchar\limits{%
    \mint@l{#1}%
  }{%
    \@ifnextchar\nolimits{%
      \mint@l{#1}%
    }{%
      \@ifnextchar\displaylimits{%
        \mint@l{#1}%
      }{%
        \mint@s{#2}{#1}%
      }%
    }%
  }%
}
\newcommand*{\mint@s}[2]{%
  % #1: limits
  % #2: overlay symbol
  \@ifnextchar_{%
    \mint@sub{#1}{#2}%
  }{%
    \@ifnextchar^{%
      \mint@sup{#1}{#2}%
    }{%
      \mint@{#1}{#2}{}{}%
    }%
  }%
}
\def\mint@sub#1#2_#3{%
  \@ifnextchar^{%
    \mint@sub@sup{#1}{#2}{#3}%
  }{%
    \mint@{#1}{#2}{#3}{}%
  }%
}
\def\mint@sup#1#2^#3{%
  \@ifnextchar_{%
    \mint@sup@sub{#1}{#2}{#3}%
  }{%
    \mint@{#1}{#2}{}{#3}%
  }%
}
\def\mint@sub@sup#1#2#3^#4{%
  \mint@{#1}{#2}{#3}{#4}%
}
\def\mint@sup@sub#1#2#3_#4{%
  \mint@{#1}{#2}{#4}{#3}%
}
\newcommand*{\mint@}[4]{%
  % #1: \limits, \nolimits, \displaylimits
  % #2: overlay symbol: -, =, ...
  % #3: subscript
  % #4: superscript
  \mathop{}%
  \mkern-\thinmuskip
  \mathchoice{%
    \mint@@{#1}{#2}{#3}{#4}%
        \displaystyle\textstyle\scriptstyle
  }{%
    \mint@@{#1}{#2}{#3}{#4}%
        \textstyle\scriptstyle\scriptstyle
  }{%
    \mint@@{#1}{#2}{#3}{#4}%
        \scriptstyle\scriptscriptstyle\scriptscriptstyle
  }{%
    \mint@@{#1}{#2}{#3}{#4}%
        \scriptscriptstyle\scriptscriptstyle\scriptscriptstyle
  }%
  \mkern-\thinmuskip
  \int#1%
  \ifx\\#3\\\else_{#3}\fi
  \ifx\\#4\\\else^{#4}\fi  
}
\newcommand*{\mint@@}[7]{%
  % #1: limits
  % #2: overlay symbol
  % #3: subscript
  % #4: superscript
  % #5: math style
  % #6: math style for overlay symbol
  % #7: math style for subscript/superscript
  \begingroup
    \sbox0{$#5\int\m@th$}%
    \sbox2{$#5\int_{}\m@th$}%
    \dimen2=\wd0 %
    % => \dimen2 = width of \int
    \let\mint@limits=#1\relax
    \ifx\mint@limits\relax
      \sbox4{$#5\int_{\kern1sp}^{\kern1sp}\m@th$}%
      \ifdim\wd4>\wd2 %
        \let\mint@limits=\nolimits
      \else
        \let\mint@limits=\limits
      \fi
    \fi
    \ifx\mint@limits\displaylimits
      \ifx#5\displaystyle
        \let\mint@limits=\limits
      \fi
    \fi
    \ifx\mint@limits\limits
      \sbox0{$#7#3\m@th$}%
      \sbox2{$#7#4\m@th$}%
      \ifdim\wd0>\dimen2 %
        \dimen2=\wd0 %
      \fi
      \ifdim\wd2>\dimen2 %
        \dimen2=\wd2 %
      \fi
    \fi
    \rlap{%
      $#5%
        \vcenter{%
          \hbox to\dimen2{%
            \hss
            $#6{#2}\m@th$%
            \hss
          }%
        }%
      $%
    }%
  \endgroup
}
\newcommand{\avg}{\mint{-}}
\makeatother

% ============================================================
% Comments and drafting helpers
% ============================================================

%Comments
\newcommand{\com}[1]{\vspace{5 mm}\par \noindent
\marginpar{\textsc{Comment}} \framebox{\begin{minipage}[c]{0.95
\textwidth} \tt #1 \end{minipage}}\vspace{5 mm}\par}

% ============================================================
% Arrows
% ============================================================

\newcommand{\hooktwoheadrightarrow}{%
  \hookrightarrow\mathrel{\mspace{-15mu}}\rightarrow
}

\makeatletter
\newcommand{\xhooktwoheadrightarrow}[2][]{%
  \lhook\joinrel
  \ext@arrow 0359\rightarrowfill@ {#1}{#2}%
  \mathrel{\mspace{-15mu}}\rightarrow
}
\makeatother

\makeatletter
\newcommand*{\da@rightarrow}{\mathchar"0\hexnumber@\symAMSa 4B }
\newcommand*{\da@leftarrow}{\mathchar"0\hexnumber@\symAMSa 4C }
\newcommand*{\xdashrightarrow}[2][]{%
  \mathrel{%
    \mathpalette{\da@xarrow{#1}{#2}{}\da@rightarrow{\,}{}}{}%
  }%
}
\newcommand{\xdashleftarrow}[2][]{%
  \mathrel{%
    \mathpalette{\da@xarrow{#1}{#2}\da@leftarrow{}{}{\,}}{}%
  }%
}
\newcommand*{\da@xarrow}[7]{%
  % #1: below
  % #2: above
  % #3: arrow left
  % #4: arrow right
  % #5: space left 
  % #6: space right
  % #7: math style 
  \sbox0{$\ifx#7\scriptstyle\scriptscriptstyle\else\scriptstyle\fi#5#1#6\m@th$}%
  \sbox2{$\ifx#7\scriptstyle\scriptscriptstyle\else\scriptstyle\fi#5#2#6\m@th$}%
  \sbox4{$#7\dabar@\m@th$}%
  \dimen@=\wd0 %
  \ifdim\wd2 >\dimen@
    \dimen@=\wd2 %   
  \fi
  \count@=2 %
  \def\da@bars{\dabar@\dabar@}%
  \@whiledim\count@\wd4<\dimen@\do{%
    \advance\count@\@ne
    \expandafter\def\expandafter\da@bars\expandafter{%
      \da@bars
      \dabar@ 
    }%
  }%  
  \mathrel{#3}%
  \mathrel{%   
    \mathop{\da@bars}\limits
    \ifx\\#1\\%
    \else
      _{\copy0}%
    \fi
    \ifx\\#2\\%
    \else
      ^{\copy2}%
    \fi
  }%   
  \mathrel{#4}%
}
\makeatother

% ============================================================
% Relation and delimiter spacing
% ============================================================

\renewcommand{\geq}{\geqslant}
\renewcommand{\leq}{\leqslant}
\newcommand{\midd}{\hspace{4pt}\middle|\hspace{4pt}}

% ============================================================
% Left Right Updated Deliminators
% ============================================================

\makeatletter
\NewDocumentCommand\Left{O{}}{%
    \IfValueTF{#1}{\notblank{#1}{\mathopen\bBigg@{#1}}{\left}}{\left}}
\NewDocumentCommand\Right{O{}}{%
    \IfValueTF{#1}{\notblank{#1}{\mathopen\bBigg@{#1}}{\right}}{\right}}
\makeatother

% ============================================================
% Restrictions
% ============================================================

\newcommand{\littletaller}{\mathchoice{\vphantom{\big|}}{}{}{}}
\newcommand\restr[2]{{% we make the whole thing an ordinary symbol
  \left.\kern-\nulldelimiterspace % automatically resize the bar with \right
  #1 % the function
  \littletaller % pretend it's a little taller at normal size
  \right|_{#2} % this is the delimiter
  }}

% ============================================================
% Chapter title formatting
% ============================================================

\newcommand{\chnum}{\titleformat
{\chapter} % command
[display] % shape
{\centering} % format
{\Huge \color{black} \shadowbox{\thechapter}} % label
{-0.5em} % sep (space between the number and title)
{\LARGE \color{black} \underline} % before-code
}

\newcommand{\chunnum}{\titleformat
{\chapter} % command
[display] % shape
{} % format
{} % label
{0em} % sep
{ \begin{flushright} \begin{tabular}{r}  \Huge \color{black}
} % before-code
[
\end{tabular} \end{flushright} \normalsize
] % after-code
}



\begin{document}
\begin{center}
    {\Large Math 502 \\[0.1in]Homework \#1}\\[.175in]
    {Name:} {\underline{Gianluca Crescenzo\hspace*{2in}}}\\[0.15in]
    \end{center}
    \vspace{4pt}
%%%%%%%%%%%%%%%%%%%%%%%%%%%%%%%%%%%%%%%%%%%%%%%%%%
\begin{exercise}[manual-num={0}]
In a short (less than half a page) paragraph, tell me about yourself and what you are hoping to get out
of this course. For example, is there a numerical component to your thesis research? Is there a certain
topic you were hoping to cover this quarter?
\end{exercise}

\begin{answer}
My name is Gianluca Crescenzo. I'm in the mathematics masters program, and I did my undergrad at Occidental College in Los Angeles. Besides math I like rock climbing, hiking, and playing videogames. I'm taking this class because I don't have much experience or proficiency in applied math. The only applied courses I've taken during my undergrad were partial differential equations, numerical analysis, and computational mathematics.
\end{answer}
%%%%%%%%%%%%%%%%%%%%%%%%%%%%%%%%%%%%%%%%%%%%%%%%%%
\begin{exercise}[manual-num={A.1}]
Let $e = (2,-5,3,4)$. Calculate the following quantities:
    \begin{equation*}
    \begin{split}
        \text{(a)}\h3\lnorm e \rnorm_\infty\h9\text{(b)}\h3\lnorm e \rnorm_1\h9\text{(c)}\h3\lnorm e \rnorm_2
    \end{split}
    \end{equation*}
\end{exercise}

\begin{solution}
    \begin{equation*}
    \begin{split}
        \lnorm e \rnorm_\infty & = \max\bigl\{|2|,|-5|,|3|,|4|\bigr\} = 5.\\
        \lnorm e \rnorm_1 & = |2| + |-5| + |3| + |4| = 14.\\
        \lnorm e \rnorm_2 & = \sqrt{2^2 +(-5)^2 + 3^2 + 4^2} = \sqrt{4 + 25 + 9 + 16} = \sqrt{54}.\\
    \end{split}
    \end{equation*}
\end{solution}
%%%%%%%%%%%%%%%%%%%%%%%%%%%%%%%%%%%%%%%%%%%%%%%%%%
\begin{exercise}[manual-num={A.2}]
    Let $A = \bmat 1 & 2 & 3 & 4 \\ 5 & 6 & 7 & 8 \\ 9 & 10 & 11 & 12 \\ 13 & 14 & 15 & 16\emat$. Calculate the following:
        \begin{equation*}
        \begin{split}
            \text{(a)}\h3\lnorm A \rnorm_\infty\h9\text{(b)}\h3\lnorm A \rnorm_1\h9\text{(c)}\h3\lnorm A \rnorm_2
        \end{split}
        \end{equation*}
\end{exercise}

\begin{solution}
    Below is the \textsc{MATLAB} code used to solve this exercise:
    \begin{lstlisting}[style=MatlabEditorGrey]
A = [1 2 3 4;5 6 7 8;9 10 11 12;13 14 15 16];
L_infty = max(sum(abs(A),2))
L1 = max(sum(abs(A),1))
L2 = sqrt(max(eig(A'*A)))
    \end{lstlisting}

    \begin{lstlisting}[style=Matlab-console]
>> untitled

L_infty =

    58

L1 =

    40

L2 =

38.6227

>>
    \end{lstlisting}
\end{solution}
%%%%%%%%%%%%%%%%%%%%%%%%%%%%%%%%%%%%%%%%%%%%%%%%%%
\begin{exercise}[manual-num={A.3}]
    Let $u(x) = x^2 - x$ be a function defined on $0 \leq x \leq 1$. Calculate the following quantities:
    \begin{equation*}
    \begin{split}
        \text{(a)}\h3\lnorm u \rnorm_\infty\h9\text{(b)}\h3\lnorm u \rnorm_1\h9\text{(c)}\h3\lnorm u \rnorm_2
    \end{split}
    \end{equation*}
\end{exercise}

\begin{solution}
    The infinity norm of $u(x) = x^2 - x$ is $\lnorm u \rnorm_\infty = \sup\{|u(x)| \mid x \in [0,1]\} = 1/2$. The one norm and two norm can be seen from the following MATLAB code:
\end{solution}

    \begin{lstlisting}[style=MatlabEditorGrey]
syms x
u = x^2 - x;
L1_exact = int(abs(u), x, 0, 1)
L2_exact = sqrt(int(abs(u)^2, x, 0, 1))
    \end{lstlisting}
    \begin{lstlisting}[style=Matlab-console]
>> untitled
 
L1_exact =
 
1/6
 
 
L2_exact =
 
30^(1/2)/30

>>
    \end{lstlisting}
%%%%%%%%%%%%%%%%%%%%%%%%%%%%%%%%%%%%%%%%%%%%%%%%%%
\begin{exercise}[manual-num={A.4}]
    Let $U$ be the grid function defined by $U_i = x_i^2 - x_i$, where $x_1 = 0.2$, $x_2 = 0.4$, $x_3 = 0.6$, $x_4 = 0.8$, $x_5 = 1.0$. Calculate the following quantities:
    \begin{equation*}
    \begin{split}
        \text{(a)}\h3\lnorm U \rnorm_\infty\h9\text{(b)}\h3\lnorm U \rnorm_1\h9\text{(c)}\h3\lnorm U \rnorm_2
    \end{split}
    \end{equation*}
\end{exercise}
\begin{solution}
    Below is the \textsc{MATLAB} code used to solve this exercise:
    \begin{lstlisting}[style=MatlabEditorGrey]
syms x
u = x^2 - x;
h = [0.2 0.4 0.6 0.8 1.0];
U = double(subs(u, x, h)) ;
Uinf = max(abs(U))
U1 = (0.2) * sum(abs(U))
U2 = sqrt((0.2) * sum(abs(U).^2))
    \end{lstlisting}
    \begin{lstlisting}[style=Matlab-console]
>> untitled

Uinf =

    0.2400


U1 =

    0.1600


U2 =

    0.1824

>>
    \end{lstlisting}
\end{solution}
%%%%%%%%%%%%%%%%%%%%%%%%%%%%%%%%%%%%%%%%%%%%%%%%%%
\begin{exercise}[manual-num={1.2}]
    \phantom{a}
\begin{enumerate}[label = (\arabic*),itemsep=1pt,topsep=3pt]
    \item Use the method of undetermined coefficients to set up the $5 \times 5$ Vandermonde system that would determine a fourth-order accurate finite difference approximation to $u''(x)$ based on $5$ equally spaced points:
        \begin{equation*}
        \begin{split}
            u''(x) = c_{-2}u(x-2h) + c_{-1}u(x-h) + c_0u(x) + c_{1}u(x+h)  c_{2}u(x+2h)
        \end{split}
        \end{equation*}
    \item Compute the coefficients using the \textsc{MATLAB} code \texttt{fdstencil.m}, and check that they satisfy the system you determined in part (1).
    \item Test this finite difference approximate $u''(1)$ for $u(x) = \sin(2x)$ with values of $h$ from the array \texttt{hvals = logspace(-1,-4,13)}. Make a table of the error vs $h$ for several values of $h$ and compare against the predicted error from the leading term of the expression printed by \texttt{fdstencil.m}. Also produce a log-log plot of the absolute value of the error vs $h$. 
\end{enumerate}
\end{exercise}

\begin{solution}
    \phantom{a}
\begin{enumerate}[label = (\arabic*),itemsep=1pt,topsep=3pt]

\item Taylor expanding $u(x-2h)$, $u(x-h)$, $u(x+h)$, and $u(x+2h)$ gives:
    \begin{equation*}
\begin{split}
    u''(x)
    &= c_{-2} \Biggl( u(x) + (-2h)u'(x) + \frac{(-2h)^2}{2!}u''(x) + \frac{(-2h)^3}{3!}u'''(x) + \frac{(-2h)^4}{4!}u''''(x) + O(h^5) \Biggr) \\
    &+ c_{-1} \Biggl( u(x) + (-h)u'(x) + \frac{(-h)^2}{2!}u''(x) + \frac{(-h)^3}{3!}u'''(x) + \frac{(-2h)^4}{4!}u''''(x) + O(h^5) \Biggr) \\
    &+ c_{0} u(x) \\
    &+ c_{1} \Biggl( u(x) + (h)u'(x) + \frac{(h)^2}{2!}u''(x) + \frac{(h)^3}{3!}u'''(x) + \frac{(-2h)^4}{4!}u''''(x) + O(h^5) \Biggr) \\
    &+ c_{2} \Biggl( u(x) + (2h)u'(x) + \frac{(2h)^2}{2!}u''(x) + \frac{(2h)^3}{3!}u'''(x) + \frac{(-2h)^4}{4!}u''''(x) + O(h^5) \Biggr).
\end{split}
\end{equation*}
Expanding each of the constants and factoring out the respective derivatives of $u(x)$ gives:
    \begin{equation*}
    \begin{split}
        u''(x) 
        & = (c_{-2}+c_{-1}+c_{0}+c_{1}+c_{2})u(x) \\
        &+ \bigl(c_{-2}(-2h) + c_{-1}(-h) + c_{1}(h) + c_{2}(2h)\bigr)u'(x) \\
        &+ \left(c_{-2}\frac{(-2h)^2}{2!} + c_{-1}\frac{(-h)^2}{2!}  + c_{1}\frac{(h)^2}{2!} + c_{2}\frac{(2h)^2}{2!}\right)u''(x) \\
        &+ \left(c_{-2}\frac{(-2h)^3}{3!} + c_{-1}\frac{(-h)^3}{3!}  + c_{1}\frac{(h)^3}{3!} + c_{2}\frac{(2h)^3}{3!}\right)u'''(x) \\
        &+ \left(c_{-2}\frac{(-2h)^4}{4!} + c_{-1}\frac{(-h)^4}{4!}  + c_{1}\frac{(h)^4}{4!} + c_{2}\frac{(2h)^4}{4!}\right)u'''(x) \\
        &+ O(h^5).
    \end{split}
    \end{equation*}
We can rewrite this as the following system of equations:
    \begin{equation*}
    \begin{split}
        c_{-2}+c_{-1}+c_{0}+c_{1}+c_{2} &= 0 \\
        c_{-2}(-2h) + c_{-1}(-h) + c_{1}(h) + c_{2}(2h) &= 0 \\
        c_{-2}\frac{(-2h)^2}{2!} + c_{-1}\frac{(-h)^2}{2!}  + c_{1}\frac{(h)^2}{2!} + c_{2}\frac{(2h)^2}{2!} & = 1 \\
        c_{-2}\frac{(-2h)^3}{3!} + c_{-1}\frac{(-h)^3}{3!}  + c_{1}\frac{(h)^3}{3!} + c_{2}\frac{(2h)^3}{3!} & = 0 \\
        c_{-2}\frac{(-2h)^4}{4!} + c_{-1}\frac{(-h)^4}{4!}  + c_{1}\frac{(h)^4}{4!} + c_{2}\frac{(2h)^4}{4!} & = 0.
    \end{split}
    \end{equation*}
Hence we've obtained the following matrix equation:
    \begin{equation*}
    \begin{split}
        \bmat
        1 & 1 & 1 & 1 & 1 \\
        -2h & -h & 0 & h & 2h \\
        \frac{(-2h)^2}{2!} & \frac{(-h)^2}{2!} & 0 & \frac{h^2}{2!} & \frac{(2h)^2}{2!}\\
        \frac{(-2h)^3}{3!} & \frac{(-h)^3}{3!} & 0 & \frac{h^3}{3!} & \frac{(2h)^3}{3!}\\
        \frac{(-2h)^4}{4!} & \frac{(-h)^4}{4!} & 0 & \frac{h^4}{4!} & \frac{(2h)^4}{4!}\\
        \emat
        \bmat c_{-2} \\ c_{-1} \\ c_{0} \\ c_{1} \\ c_{2} \emat  
        = \bmat 0 \\ 0 \\ 1 \\ 0\\ 0 \emat .
    \end{split}
    \end{equation*}
Factoring an $h^2$ out of the third row and simplifying yields:
    \begin{equation*}
    \begin{split}
        \bmat
        1 & 1 & 1 & 1 & 1 \\
        -2h & -h & 0 & h & 2h \\
        2 & \frac{1}{2} & 0 & \frac{1}{2} & 2\\
        -\frac{4h^3}{3} & -\frac{h^3}{6} & 0 & \frac{h^3}{6} & \frac{4h^3}{3}\\
        \frac{2h^4}{3} & \frac{h^4}{24} & 0 & \frac{h^4}{24} & \frac{2h^4}{3}\\
        \emat
        \bmat c_{-2} \\ c_{-1} \\ c_{0} \\ c_{1} \\ c_{2} \emat  
        = \bmat 0 \\ 0 \\ \frac{1}{h^2} \\ 0\\ 0 \emat .
    \end{split}
    \end{equation*}

\item The following code row reduces the the augmented matrix obtained from the above matrix equation, as well as prints the output of \texttt{fdstencil(2,-2:2)}.
    \begin{lstlisting}[style=MatlabEditorGrey]
syms h

A = [      1,        1, 1,        1,         1,     0;
        -2*h,       -h, 0,        h,       2*h,     0;
           2,      1/2, 0,      1/2,         2, 1/h^2;
  -(4*h^3)/3, -(h^3)/6, 0,  (h^3)/6, (4*h^3)/3,     0;
   (2*h^4)/3, (h^4)/24, 0, (h^4)/24, (2*h^4)/3,     0];

disp(rref(A))
fdstencil(2,-2:2)
    \end{lstlisting}
    \newpage
    \begin{lstlisting}[style=Matlab-console]
>> untitled
[1, 0, 0, 0, 0, -1/(12*h^2)]
[0, 1, 0, 0, 0,   4/(3*h^2)]
[0, 0, 1, 0, 0,  -5/(2*h^2)]
[0, 0, 0, 1, 0,   4/(3*h^2)]
[0, 0, 0, 0, 1, -1/(12*h^2)]
 
 
The derivative u^(2) of u at x0 is approximated by
 
 1/h^2 * [
 -8.333333333333333e-02 * u(x0-2*h) + 
  1.333333333333333e+00 * u(x0-1*h) + 
 -2.500000000000000e+00 * u(x0) + 
  1.333333333333333e+00 * u(x0+1*h) + 
 -8.333333333333333e-02 * u(x0+2*h) ] 
 
For smooth u,
 Error = 0 * h^3*u^(5) + -0.0111111 * h^4*u^(6) + ...
 
>> 
    \end{lstlisting}
    Hence the coefficients obtained by \texttt{fdstencil.m} satisfy the system we determined in part (1).
\item The following \textsc{MATLAB} code and output addresses the question:
    \begin{lstlisting}[style=MatlabEditorGrey]
% u(x) = sin(2x), approximate u''(1):
uptrue = -4*sin(2);
%hvals = [1e-1 5e-2 1e-2 5e-3 1e-3];
hvals = logspace(-1, -4, 13);
E5u = [];
% table headings:
disp(' ')
disp(' h D3u E3u')
for i=1:length(hvals)
h = hvals(i);
% approximation to u''(1) using fdcoeffF:
xpts = 1 + h*(-2:2)';
D5u = fdcoeffF(2,1,xpts) * sin(2*xpts) ;
% error:
E5u(i) = D5u - uptrue;
% print line of table:
disp(sprintf('%13.4e %13.4e %13.4e',...
h,D5u,E5u(i)))
end
% plot errors:
clf
loglog(hvals,abs(E5u),'o-')
axis([5e-7 .2 1e-12 1])
    \end{lstlisting}
    \begin{lstlisting}[style=Matlab-console]
>> untitled
 
   h             D3u            E3u
   1.0000e-01   -3.6371e+00    6.4431e-05
   5.6234e-02   -3.6372e+00    6.4588e-06
   3.1623e-02   -3.6372e+00    6.4638e-07
   1.7783e-02   -3.6372e+00    6.4655e-08
   1.0000e-02   -3.6372e+00    6.4670e-09
   5.6234e-03   -3.6372e+00    6.4623e-10
   3.1623e-03   -3.6372e+00    6.4150e-11
   1.7783e-03   -3.6372e+00   -6.6817e-11
   1.0000e-03   -3.6372e+00   -1.8323e-10
   5.6234e-04   -3.6372e+00   -1.8130e-09
   3.1623e-04   -3.6372e+00    4.2406e-09
   1.7783e-04   -3.6372e+00    6.1032e-09
   1.0000e-04   -3.6372e+00   -3.8600e-08
>> 
    \end{lstlisting}

\newpage
\begin{center}
\scalebox{0.7}{\definecolor{mycolor1}{rgb}{0.14900,0.54900,0.86600}
\definecolor{mycolor2}{rgb}{0.07059,0.07059,0.07059}
\begin{tikzpicture}
\begin{axis}[
width=7.887in,
height=4.992in,
at={(1.323in,0.674in)},
scale only axis,
xmode=log,
xmin=5e-07,
xmax=0.2,
xminorticks=true,
ymode=log,
ymin=1e-12,
ymax=1,
yminorticks=true,
axis background/.style={fill=white}
]
\addplot [color=mycolor1, mark=o, mark options={solid, mycolor1}, forget plot]
  table[row sep=crcr]{
0.1	6.4430648182956e-05\\
0.0562341325190349	6.45881661665015e-06\\
0.0316227766016838	6.46380871494046e-07\\
0.0177827941003892	6.46546451932295e-08\\
0.01	6.4669931632011e-09\\
0.00562341325190349	6.46227071854355e-10\\
0.00316227766016838	6.41504627196809e-11\\
0.00177827941003892	6.68167743356207e-11\\
0.001	1.83232096162556e-10\\
0.000562341325190349	1.81304660173964e-09\\
0.000316227766016838	4.24055013326097e-09\\
0.000177827941003892	6.10319528249192e-09\\
0.0001	3.8600288299051e-08\\
};
\end{axis}
\end{tikzpicture}}
\end{center}


\end{enumerate}
\end{solution}
%%%%%%%%%%%%%%%%%%%%%%%%%%%%%%%%%%%%%%%%%%%%%%%%%%

\end{document}